\chapter{Pulmonary Embolism}

\section*{Patient focus}
This chapter explains the examination or clinical problem in accessible language. Patients should use it to prepare questions and understand the purpose, experience, benefits, and limitations of imaging; local instructions and the official report take priority.


\label{ch:ch73_pulmonary_embolism}

A pulmonary embolism (PE) is a blood clot that travels from the legs or pelvis to the lungs, blocking a pulmonary artery. PE is common, potentially fatal, and frequently missed. Imaging plays a central role in confirming or excluding the diagnosis, assessing severity, and guiding treatment decisions.

---

\section{1. Patient Perspective: What the PE Scan Involves}

If your doctor suspects a blood clot in your lungs, you will likely have a CT scan called a \textbf{CT pulmonary angiogram} (CTPA). This scan uses X-rays and injected contrast dye to produce detailed pictures of the blood vessels in your lungs.

\begin{itemize}
  \item You will lie on a scanner table that slides through a doughnut-shaped machine.
  \item An IV cannula is placed, usually in the arm. Contrast dye is injected, which may cause a brief warm flush or metallic taste.
  \item The scan itself takes only a few seconds. You will be asked to hold your breath briefly.
  \item The total visit lasts about 30 minutes.
  \item Iodinated contrast can cause an immediate reaction and requires individualized kidney-risk assessment. Kidney blood tests are obtained when indicated by your history and local protocol; they are not required routinely for every patient.
\end{itemize}

If you are pregnant or have a contrast allergy, alternative tests such as a ventilation-perfusion (V/Q) scan may be arranged instead.

---

\section{2. The Clinical Question and Pretest Probability}

PE should be considered in any patient with sudden-onset breathlessness, pleuritic chest pain, haemoptysis, unexplained tachycardia, or syncope. The \textbf{Wells score} stratifies pretest probability:

\begin{table}[h]
\centering
\begin{tabular}{l l}
\hline
\textbf{Criterion} & \textbf{Points} \\
\hline
Clinical signs of DVT & +3 \\
Haemoptysis & +1 \\
Heart rate $>$ 100 bpm & +1.5 \\
Immobilisation or recent surgery & +1.5 \\
Malignancy & +1 \\
PE is \#1 diagnosis or equally likely & +3 \\
Previous DVT/PE & +1.5 \\
\hline
\multicolumn{2}{l}{Three-level Wells score: low $\leq1$; intermediate 2--6; high $>6$. Some pathways use the two-level cutoff: PE unlikely $\leq4$, PE likely $>4$.} \\
\hline
\end{tabular}
\caption{Wells Score for Pulmonary Embolism}
\end{table}

Use a validated pathway consistently. PERC may avoid testing only when the clinician has already judged PE probability to be very low and all PERC criteria are negative. In low or intermediate probability, a negative D-dimer below the pathway-specific (including age-adjusted or YEARS) threshold can avoid imaging. High clinical probability warrants imaging; a negative D-dimer alone is not sufficient to rule out PE.

---

\section{3. CTPA: Protocol and Technical Considerations}

\begin{itemize}
  \item \textbf{Timing}: Contrast bolus is timed to opacify the pulmonary arteries (bolus tracking at the main pulmonary artery, trigger $\sim$100 HU).
  \item \textbf{Phase}: Arterial phase acquisition, ideally without cardiac motion artifact.
  \item \textbf{Radiation}: Dose varies with scanner, patient, and protocol; a single effective-dose estimate is not a patient-specific risk measure.
  \item \textbf{Contrast}: 80--100 mL of iodinated contrast at 4--5 mL/s, followed by saline chase.
\end{itemize}

---

\section{4. PE Imaging Findings}

\subsection{Central vs Peripheral Embolism}
\begin{itemize}
  \item \textbf{Central PE}: Emboli in the main or lobar arteries; a saddle embolus straddles the main pulmonary artery bifurcation. Clot location alone does not define clinical risk or hemodynamic severity.
  \item \textbf{Segmental PE}: Emboli within segmental pulmonary arteries. Often present with pleuritic pain and haemoptysis.
  \item \textbf{Subsegmental PE}: Emboli confined to subsegmental branches. These are harder to detect and their clinical significance is debated, particularly in patients without DVT and with low pretest probability.
\end{itemize}

\subsection{Right Heart Strain}
\begin{itemize}
  \item \textbf{RV/LV ratio} $>$ 1.0 on axial CT (measured at the level of the papillary muscles) indicates right ventricular dilatation.
  \item \textbf{Bowing of the interventricular septum} toward the left ventricle.
  \item \textbf{Contrast reflux} into the IVC and hepatic veins suggests elevated right heart pressures.
\end{itemize}

\subsection{Infarction}
Peripheral wedge-shaped opacities with a broad pleural base (Hampton's hump) can support pulmonary infarction in the appropriate setting, but are neither sensitive nor specific for PE.

\section{5. PE Severity Index and Imaging Role}

The \textbf{Pulmonary Embolism Severity Index (PESI)} and its simplified version (sPESI) guide risk stratification. Imaging contributes as follows:

\begin{itemize}
  \item Use current risk terminology and clinical classification. Hemodynamic instability defines high-risk PE; central clot location alone does not.
  \item In normotensive patients, integrate RV imaging findings, biomarkers, clinical scores, comorbidities, and bleeding risk. RV strain does not by itself mandate thrombolysis or catheter intervention.
  \item Stable patients without concerning clinical or imaging features may be candidates for anticoagulation and selected outpatient care after validated assessment.
\end{itemize}

\section{6. Differential Diagnosis}

CTPA also identifies alternative diagnoses that may mimic PE:
\begin{itemize}
  \item Pneumonia or pulmonary infarction
  \item Pleurisy or pleural effusion
  \item Pericardial effusion or pericarditis
  \item Aortic dissection (if imaging extends to the aorta)
  \item Anxiety or hyperventilation
  \item Pneumothorax
\end{itemize}

\begin{pearlbox}[CTPA Pitfall]
Inadequate pulmonary arterial opacification or motion can limit assessment, especially of segmental and subsegmental vessels. Judge adequacy from the acquisition protocol and visible opacification at the vessel levels being assessed; no single attenuation cutoff is universally valid. State limitations and recommend follow-up only when clinically appropriate.
\end{pearlbox}

\section{7. V/Q Scan: Role and Limitations}

Ventilation-perfusion scintigraphy is an alternative when CTPA is contraindicated:
\begin{itemize}
  \item May be selected in pregnancy, particularly when the chest radiograph is normal and the local protocol supports it; the choice between V/Q and CTPA depends on patient factors, availability, and the expected diagnostic yield and dose.
  \item A \textbf{normal perfusion or normal V/Q scan} effectively excludes clinically significant acute PE in the appropriate diagnostic context.
  \item A low-probability result is not, by itself, a rule-out when clinical pretest probability is high; interpret the scan with clinical probability and the reporting scheme.
  \item A \textbf{high-probability scan} (mismatched defects) is strongly suggestive of PE.
  \item Planar scans can be nondiagnostic, especially with abnormal chest imaging; V/Q SPECT or SPECT/CT may reduce nondiagnostic interpretations where available.
  \item Less useful in patients with pre-existing lung disease (COPD, prior PE).
\end{itemize}

---

\section{8. Structured Reporting Example}

\begin{testbox}[Indications:]
52-year-old female, 2 weeks post knee replacement, presenting with sudden-onset dyspnoea and pleuritic right-sided chest pain. Wells score 6 (high probability). D-dimer 4,200 ng/mL.
\end{testbox}
\begin{refbox}[Findings:]
\textbf{Pulmonary Vasculature}: Filling defects are identified within the right lower lobe segmental pulmonary arteries and left lower lobe segmental arteries, consistent with bilateral segmental pulmonary emboli. The main, right, and left pulmonary arteries are clear. No saddle embolus.\\
\textbf{Right Heart}: The RV/LV ratio measures 1.3 (axial), with mild bowing of the interventricular septum. Mild contrast reflux into the IVC is noted.\\
\textbf{Lungs}: Small right-sided pleural effusion. No consolidation, ground-glass opacity, or pneumothorax.\\
\textbf{Other}: No mediastinal lymphadenopathy. Bones are intact.\\
\textbf{IVC Filter}: None in situ.
\end{refbox}
\begin{pearlbox}[Impression:]
1. Bilateral segmental pulmonary emboli with CT evidence of right ventricular enlargement/strain (RV/LV ratio 1.3).
2. Small right pleural effusion.
3. Hemodynamic PE risk category requires clinical assessment, including blood pressure, biomarkers, and the patient's overall status. Correlate with echocardiography when it will affect management.
\end{pearlbox}

\section*{Selected Guidance}
\begin{itemize}
  \item American Heart Association/American College of Cardiology and collaborating societies. \href{https://www.ahajournals.org/doi/pdf/10.1161/CIR.0000000000001415}{2026 Guideline for Evaluation and Management of Acute Pulmonary Embolism in Adults}.
  \item Society of Nuclear Medicine and Molecular Imaging and collaborating societies. \href{https://doi.org/10.2967/jnmt.126.273453}{Procedure standard for ventilation--perfusion pulmonary scintigraphy}.
\end{itemize}

\clearpage
